\section*{अध्याय \ref{ch:g9:sqrt} --- वर्गमूल}

\begin{solution}{exo:g9:sqrt:1}
$\sqrt{64} = 8$; \quad $\sqrt{100} = 10$; \quad
$\sqrt{0.09} = 0.3$ (क्योंकि $0.3^2 = 0.09$); \quad
$\sqrt{\frac{1}{25}} = \frac15$; \quad
$\left(\sqrt{13}\right)^2 = 13$; \quad
$\sqrt{(-4)^2} = \sqrt{16} = 4 = \abs{-4}$.
\end{solution}

\begin{solution}{exo:g9:sqrt:2}
$\sqrt{50} = \sqrt{25 \times 2} = 5\sqrt2$; \quad
$\sqrt{45} = \sqrt{9 \times 5} = 3\sqrt5$; \quad
$\sqrt{98} = \sqrt{49 \times 2} = 7\sqrt2$; \quad
$\sqrt{300} = \sqrt{100 \times 3} = 10\sqrt3$.
\end{solution}

\begin{solution}{exo:g9:sqrt:3}
$\sqrt2 \times \sqrt{18} = \sqrt{36} = 6$.

$\dfrac{\sqrt{75}}{\sqrt3} = \sqrt{\dfrac{75}{3}} = \sqrt{25} = 5$.

$\sqrt5 \times \sqrt{20} = \sqrt{100} = 10$.

$\dfrac{\sqrt8}{\sqrt{50}} = \sqrt{\dfrac{8}{50}} = \sqrt{\dfrac{4}{25}}
= \dfrac25$.
\end{solution}

\begin{solution}{exo:g9:sqrt:4}
$x^2 = 36$: $x = 6$ या $x = -6$।

$x^2 = 11$: $x = \sqrt{11}$ या $x = -\sqrt{11}$।

$x^2 + 4 = 0$ यानी $x^2 = -4 < 0$: कोई हल नहीं।

$2x^2 = 50$: $x^2 = 25$, इसलिए $x = 5$ या $x = -5$।
\end{solution}

\begin{solution}{exo:g9:sqrt:5}
$\sqrt{20} + \sqrt{45} = 2\sqrt5 + 3\sqrt5 = 5\sqrt5$.

$3\sqrt8 - \sqrt{18} + \sqrt2 = 3 \times 2\sqrt2 - 3\sqrt2 + \sqrt2
= (6 - 3 + 1)\sqrt2 = 4\sqrt2$.
\end{solution}

\begin{solution}{exo:g9:sqrt:6}
$\left(\sqrt3 + 1\right)^2 = 3 + 2\sqrt3 + 1 = 4 + 2\sqrt3$.

$\left(2\sqrt5 - 3\right)^2 = 4 \times 5 - 2 \times 3 \times 2\sqrt5 + 9
= 29 - 12\sqrt5$.

$\left(\sqrt6 + \sqrt2\right)\left(\sqrt6 - \sqrt2\right) = 6 - 2 = 4$.
\end{solution}

\begin{solution}{exo:g9:sqrt:7}
भुजा: $\sqrt{6400} = 80$ m। विकर्ण (पाइथागोरस से):
$\sqrt{80^2 + 80^2} = 80\sqrt2 \approx 113$ m।
\end{solution}

\begin{solution}{exo:g9:sqrt:8}
$\dfrac{1}{\sqrt2} = \dfrac{1 \times \sqrt2}{\sqrt2 \times \sqrt2}
= \dfrac{\sqrt2}{2}$.

$\dfrac{6}{\sqrt3} = \dfrac{6\sqrt3}{3} = 2\sqrt3$.
\quad
$\dfrac{10}{\sqrt5} = \dfrac{10\sqrt5}{5} = 2\sqrt5$.
\end{solution}

\begin{solution}{exo:g9:sqrt:9}
\emph{1. सही}: यह \omterm{def:g2:mult:def}{गुणनफल} वाला नियम है (\cref{thm:g9:sqrt:rules})।

\emph{2. ग़लत}: $\sqrt{9 + 16} = 5$ है पर $\sqrt9 + \sqrt{16} = 7$।

\emph{3. ग़लत} ऋण $x$ के लिए: $\sqrt{(-4)^2} = 4 \neq -4$। सही कथन
$\sqrt{x^2} = \abs{x}$ है।

\emph{4. सही}: $\left(3\sqrt2\right)^2 = 9 \times 2 = 18$।
\end{solution}

\begin{solution}{exo:g9:sqrt:10}
$\left(2 + \sqrt3\right)^2 = 4 + 4\sqrt3 + 3 = 7 + 4\sqrt3$। इसलिए
$x = \sqrt{7 + 4\sqrt3} = \sqrt{\left(2+\sqrt3\right)^2} = 2 + \sqrt3$
(यह अऋणात्मक संख्या है, इसलिए मूल सिर्फ़ वर्ग हटा देता है)।

इसी तरह $\left(\sqrt5 - 2\right)^2 = 5 - 4\sqrt5 + 4 = 9 - 4\sqrt5$,
और $\sqrt5 - 2 > 0$ है, इसलिए
$\sqrt{9 - 4\sqrt5} = \sqrt5 - 2$ ($2 - \sqrt5$ नहीं, जो ऋण
है!)।
\end{solution}

\begin{solution}{pb:g9:sqrt:1}
\textbf{1.} $1.4^2 = 1.96 < 2 < 2.25 = 1.5^2$;
$1.41^2 = 1.9881 < 2 < 2.0164 = 1.42^2$। तीन दशमलव:
$1.414^2 = 1.999396 < 2 < 2.002225 = 1.415^2$, इसलिए
$1.414 < \sqrt2 < 1.415$।

\textbf{2.} $\sqrt2 = \frac ab$ का वर्ग करने पर
$2 = \frac{a^2}{b^2}$, यानी $a^2 = 2b^2$: संख्या $a^2$ किसी पूर्ण
संख्या की दुगनी है --- \emph{सम}।

\textbf{3.} अगर $a$ विषम होता, तो $a^2$ भी विषम होता
(\cref{pb:g8:pythagoras:1}); और चूँकि $a^2$ सम है, इसलिए $a$ सम है:
$a = 2k$। रखकर देखो: $4k^2 = 2b^2$, यानी $b^2 = 2k^2$ सम है, और उसी
सम-विषम तथ्य से $b$ भी सम है।

\textbf{4.} $a$ और $b$ दोनों का सम होना यानी दोनों का $2$ से \omterm{def:g6:wholes:divisible}{विभाज्य}
होना --- पर $\frac ab$ सबसे सरल रूप में थी, जिसमें कोई साझा \omterm{def:g9:arith:divisor}{भाजक} था
ही नहीं। विरोधाभास: मानी हुई \omterm{def:g9:fractions:fraction}{भिन्न} हो ही नहीं सकती।
\emph{प्रमेय: $\sqrt2$ अपरिमेय है।}

\textbf{5.} सारा विरोधाभास ``सबसे सरल रूप'' में ही बसता है: उसके
बिना ``$a$ और $b$ दोनों सम'' किसी बात को नहीं काटता। सबसे सरल रूप
मान लेना जायज़ है क्योंकि हर \omterm{def:g9:fractions:fraction}{भिन्न} का सबसे सरल रूप होता \emph{है}
(\cref{met:g9:arith:simplify}: \omterm{def:g9:arith:gcd}{महत्तम समापवर्तक} से भाग दो) --- सो
अगर $\sqrt2$ कोई \omterm{def:g9:fractions:fraction}{भिन्न} होता भी, तो वह सबसे सरल रूप वाली \omterm{def:g9:fractions:fraction}{भिन्न} होता,
और वैसी \omterm{def:g9:fractions:fraction}{भिन्न} नामुमकिन है।

\textbf{6.} अभाज्य गुणनखंडनों में वर्गों के \omterm{def:g8:powers:def}{घातांक} सम होते हैं
(\cref{pb:g9:arith:1})। $a^2 = 3b^2$ में $3$ का \omterm{def:g8:powers:def}{घातांक} बाईं ओर सम है
पर दाईं ओर विषम ($b^2$ से आया सम, और उसमें एक और)। किसी संख्या के दो
गुणनखंडन नहीं होते (\cref{thm:g9:arith:factorization}): विरोधाभास,
और $\sqrt3$ अपरिमेय है।

\textbf{7.} $a^2 = n b^2$ में यह तर्क ठीक तभी किसी \omterm{def:g9:arith:prime}{अभाज्य संख्या} का
उलटा सम-विषम पकड़ता है जब कोई \omterm{def:g9:arith:prime}{अभाज्य संख्या} $n$ में \emph{विषम}
\omterm{def:g8:powers:def}{घातांक} के साथ आती हो। अगर $n$ के सारे \omterm{def:g8:powers:def}{घातांक} सम हों, तो $n$ पूर्ण
वर्ग है ($n = m^2$, और $\sqrt n = m$ पूर्ण संख्या)। पूरा नतीजा: हर
उस पूर्ण संख्या $n$ के लिए $\sqrt n$ अपरिमेय है जो पूर्ण वर्ग नहीं
है --- $\sqrt2, \sqrt3, \sqrt5, \sqrt6, \sqrt7, \sqrt8, \sqrt{10},
\dots$

\textbf{8.} $r x = q$ से, जहाँ $r \neq 0$ है:
$x = \frac qr$, यानी परिमेय संख्याओं का \omterm{def:g3:division:remainder}{भागफल}, जो परिमेय है। इसलिए
अगर $x$ अपरिमेय हो, तो $r x$ परिमेय हो ही नहीं सकता: $2\sqrt2$ और
$\frac{\sqrt2}{2}$ अपरिमेय हैं (गुणक $2$ और $\frac12$)।

\textbf{9.} अगर $1 + \sqrt2 = q$ परिमेय होता, तो $\sqrt2 = q - 1$ भी
परिमेय होता: विरोधाभास। और अगर $s = \sqrt2 + \sqrt3$ परिमेय होता, तो
$s^2 = 2 + 2\sqrt6 + 3 = 5 + 2\sqrt6$ भी परिमेय होता, जिससे
$\sqrt6 = \frac{s^2 - 5}{2}$ परिमेय बन जाता --- पर $\sqrt6$ अपरिमेय
है (सवाल 7)। इसलिए $\sqrt2 + \sqrt3$ अपरिमेय है।

\textbf{10.} \omterm{def:g1:addition:def}{योगफल}: $\sqrt2$ और $-\sqrt2$ (या $\sqrt2$ और
$1 - \sqrt2$) दोनों ही अपरिमेय हैं, पर उनके \omterm{def:g1:addition:def}{योगफल} $0$ और $1$ परिमेय
हैं। \omterm{def:g2:mult:def}{गुणनफल}: $\sqrt2 \times \sqrt2 = 2$। अंकगणित में अपरिमेयता उड़
सकती है --- इसीलिए हर हालत का अलग प्रमाण चाहिए।

\textbf{11.} $x = 1$ से: $1$ और $\frac21 = 2$ का औसत:
$x_1 = \frac32$। फिर $\frac32$ और $\frac{2}{3/2} = \frac43$ का औसत:
$x_2 = \frac12\left(\frac32 + \frac43\right)
= \frac12 \cdot \frac{17}{6} = \frac{17}{12} \approx 1.4167$।

\textbf{12.} $x_3 = \frac12\left(\frac{17}{12} +
\frac{24}{17}\right) = \frac12 \cdot
\frac{289 + 288}{204} = \frac{577}{408} = 1.4142156\ldots$
और $\sqrt2 = 1.4142135\ldots$ के मुक़ाबले: तीन क़दमों के बाद ही पाँच
दशमलव सही --- तरकीब हर चक्कर में सही दशमलवों की \omterm{def:g7:stats:frequency}{गिनती} क़रीब-क़रीब
दुगनी कर देती है।

\textbf{13.} अगर $x^2 > 2$ हो (अंदाज़ा बहुत बड़ा), तो $x > \sqrt2$
और $\frac2x < \frac{2}{\sqrt2} = \sqrt2$: साथी भुजा बहुत छोटी है।
अगर $x^2 < 2$ हो, तो असमिकाएँ उलट जाती हैं। दोनों हालतों में
$\sqrt2$, $x$ और $\frac2x$ के बीच बैठता है, इसलिए उनका औसत --- यानी
अगला अंदाज़ा --- दोनों भुजाओं में से बुरी वाली से नज़दीक होता है:
\omterm{def:g6:measure:area}{क्षेत्रफल} $2$ वाला आयत हर क़दम पर और वर्गाकार होता जाता है।

\textbf{14.} $9 - 8 = 1$; \ $289 - 288 = 1$; \
$332\,929 - 332\,928 = 1$। हेरॉन का हर अंदाज़ा $\frac ab$
$a^2 - 2b^2 = 1$ पूरा करता है: भाग I ने सिद्ध किया कि $0$ पर
पहुँचना नामुमकिन है, और ये भिन्नें उस नामुमकिन से ठीक एक इकाई चूकती
हैं --- पूर्ण संख्याएँ इससे ज़्यादा पास कभी नहीं आ सकतीं।

\textbf{15.} संख्या रेखा में भिन्नों से आगे की संख्याएँ होनी ही
चाहिए --- $\sqrt2$ जैसी लंबाइयाँ, जिनका दशमलव लेखन बिना दोहराए
हमेशा चलता रहता है: यानी \emph{वास्तविक संख्याएँ}। उनका ईमानदार
निर्माण हाई-स्कूल खंड की कहानी है, और पूरी कड़ाई के साथ
विश्वविद्यालय वाले खंडों की।

\end{solution}
